\section{The spread}\label{sec:partb}

\subsection{Target, grid and controls}

The target is the spread $\spread = \text{HH}_t - \text{TTF}_t$ in \$/MMBtu, with TTF
converted from EUR/MWh at the day's EUR/USD rate and 3.412\,MWh per MMBtu. It is
sampled weekly on Mondays from 2018, the same grid as the physical nowcasts. The
scored quantity is the $h$-week-ahead change, $\Delta_h = s_{t+7h}-s_t$, for
$h\in\{1,4\}$. The level is too persistent to score usefully, because a model that
predicts the current value performs well on it. The primary sample runs to
\noaaEnds{} (Section~\ref{sec:panel}) and contains \bWeeksCov{} complete weeks.

The control block consists of the spread's non-tanker drivers: heating and cooling
degree days for the US and north-west Europe, US and EU storage, Brent, and a winter
dummy. The price legs that construct the target are excluded as regressors. All
standard errors are Newey--West \citep{neweywest1987} with the lag fixed in advance
at $h+1$ weeks. On a weekly grid with $h=4$, consecutive targets share three weeks,
and an ordinary least-squares standard error is roughly $\sqrt{h}$ too small. Fits
are walk-forward with an expanding window and a 104-week minimum. A training row is
admissible only if its target closed strictly before the test week, with a further
one-week embargo \citep{lopezdeprado2018}.

\subsection{The controls-only null is worse than no change}

Three models form a ladder in which each is the null for the next. M0 predicts no
change. M1 is an AR(1) in the level. M2 adds the controls and is the AR(1)+controls
null that the tanker signals were pre-registered to beat. Table~\ref{tab:ladder}
shows that M2 does not beat M0. Relative to a no-change forecast, M2 scores
\bMTwoSkillHOneCov{} at one week and \bMTwoSkillHFourCov{} at four. It is worse at the
longer horizon and worse in each of 2023, 2024 and 2025 taken individually. This is
the behaviour expected of an efficient and persistent price series: the level already
incorporates the fundamentals, so adding them as regressors introduces estimation
noise. The operative null for the signals is therefore M0, a stricter bar than the
one pre-registered.

\begin{table}[t]
\centering\small
\caption{The model ladder on the primary sample. Skill is the fractional MAE
improvement over M0; a negative value means worse than a no-change forecast.}
\label{tab:ladder}
\begin{tabular}{llrrrrr}
\toprule
$h$ (wk) & model & $n$ & MAE & RMSE & 80\% cov. & skill vs M0 \\
\midrule
  1 & M0 & 272 & 2.034 & 4.465 & 81.6\% & -- \\
  1 & M1 & 272 & 2.061 & 4.571 & 81.6\% & -1.4\% \\
  1 & M2 & 272 & 2.223 & 4.630 & 82.4\% & -9.3\% \\
  4 & M0 & 266 & 3.942 & 7.203 & 72.9\% & -- \\
  4 & M1 & 266 & 4.196 & 7.585 & 70.7\% & -6.4\% \\
  4 & M2 & 266 & 4.660 & 7.437 & 69.5\% & -18.2\% \\
\bottomrule
\end{tabular}

\end{table}

\subsection{No signal survives partialling}

For each of eight pre-registered signals we residualise both the target and the
signal on M2's regressors and regress residual on residual. This is the
Frisch--Waugh--Lovell partial effect \citep{frischwaugh1933,lovell1963}. By the
theorem it equals the signal's coefficient in the full regression, and that equality
serves as the harness's correctness test. Each signal carried a sign fixed in advance
(Appendix~\ref{app:signals}). A result counted only if $|t|>1.96$, the sign matched,
and the sign held in at least two-thirds of individual years. The last condition
guards against 2022 alone driving a pooled fit.

\begin{figure}[t]
\centering
\includegraphics[width=0.9\linewidth]{fig3_fwl_scan}
\caption{Newey--West $t$-statistics on the FWL partial effect of each pre-registered
signal, after controlling for persistence, weather, storage and oil, at one and four
weeks. The dotted and dashed lines mark the two decision bars used in the paper.}
\label{fig:fwl}
\end{figure}

\begin{table}[t]
\centering\small
\caption{FWL partial effects on the primary sample. Year consistency is the share of
individual years whose partial effect has the full-sample sign.}
\label{tab:fwl}
\begin{tabular}{llrrrrr}
\toprule
 & & \multicolumn{2}{c}{$h=1$} & \multicolumn{2}{c}{$h=4$} & \\
\cmidrule(lr){3-4} \cmidrule(lr){5-6}
signal & prior & HAC $t$ & partial $R^2$ & HAC $t$ & partial $R^2$ & year cons. \\
\midrule
  gas in transit, EU-bound & $+$ & +0.18 & 0.000 & +1.40 & 0.009 & 62\% \\
  gas in transit, unresolved & none & -0.03 & 0.000 & +0.24 & 0.000 & 50\% \\
  US loading rate & $+$ & -0.32 & 0.000 & +1.14 & 0.006 & 62\% \\
  EU discharge rate & $+$ & +0.83 & 0.001 & +1.21 & 0.003 & 62\% \\
  ballast returning to US & $+$ & +0.61 & 0.001 & +1.50 & 0.008 & 62\% \\
  laden voyage age & $-$ & +0.02 & 0.000 & -1.14 & 0.004 & 62\% \\
  US load-queue wait & $-$ & +0.78 & 0.001 & +0.07 & 0.000 & 38\% \\
  net export pressure & $+$ & -0.49 & 0.000 & -1.11 & 0.003 & 75\% \\
\bottomrule
\end{tabular}

\end{table}

No signal meets the bar (Figure~\ref{fig:fwl}, Table~\ref{tab:fwl}). The largest
$|t|$ at either horizon is \bMaxTCov, and the largest \pRsq{} is \bMaxRsqCov. The
most direct mechanical hypothesis, that cargo at sea today lands in Europe in two to
three weeks, was tested as the EU-bound stock and gave $t=\bTransitTOneCov$ at one
week and $\bTransitTFourCov$ at four. At four weeks the four flow signals all have
their pre-registered sign, but all lie well inside the noise band. They are
directionally consistent with the mechanism and uniformly insignificant. No sign was
falsified, because no estimate reached significance.

\subsection{Three mechanisms a linear scan could miss}

A one-at-a-time linear scan has three identifiable blind spots. Each was tested with
a derived feature or a sample split, so each costs a single test. The three were
pre-registered together with a Bonferroni bar of $|t|>\bonferroniT$, a primary
horizon of four weeks chosen from the voyage mechanism, and a holdout
(\mechHoldLoCov{} to \mechHoldHiCov, \mechHoldNCov{} weeks) that was not examined
until a hypothesis had passed on the discovery sample (\mechDiscLoCov{} to
\mechDiscHiCov, \mechDiscNCov{} weeks). A pass required the holdout sign to match.

H2 tests the share of cargo bound for Europe in place of the level. Most of the
at-sea stock has no resolved destination, so two collinear levels with offsetting
effects could each appear insignificant while their ratio carries the signal. A
pre-registered gate voids the test if the share is mostly a decade-long trend, which
would indicate improving coverage rather than changing cargo routing. The trend
explains $\Rsq=\mechTrendRsqCov$ and the gate passes. H3 tests an interaction with EU
storage. Additional cargo arriving when storage is 90\% full has little effect,
while the same cargo at 30\% is highly price-relevant, and a linear fit averages the
two states. H4 tests a tightness split, the simplest form of a regime model. The
scan is re-run within weeks in which EU storage is below its seasonal norm (with the
norm fitted on the discovery sample only), and the prediction is that the effect is
positive in those weeks and larger than in the remaining weeks.

\begin{table}[t]
\centering\footnotesize
\caption{The three pre-registered mechanism tests on the discovery sample at $h=4$
with Bonferroni bar $|t|>\bonferroniT$. None is significant. H3 has the wrong sign,
and H4's tight-regime effect is smaller than its loose-regime effect.}
\label{tab:mechanisms}
\begin{tabular}{llrrrrr}
\toprule
hypothesis & prior & $\hat\beta$ (discovery) & HAC $t$ & year cons. & partial $R^2$ & $\hat\beta$ (holdout) \\
\midrule
  H2: EU share of at-sea stock & $+$ & $-2.64 \times 10^{0}$ & -0.89 & 71\% & 0.003 & $-2.51 \times 10^{0}$ \\
  H3: transit $\times$ EU storage & $-$ & $4.78 \times 10^{-9}$ & +0.46 & 71\% & 0.001 & $-1.49 \times 10^{-9}$ \\
  H4: transit, tight regime & $+$ & $-1.13 \times 10^{-7}$ & -0.17 & 100\% & 0.000 & $2.71 \times 10^{-7}$ \\
  H4: transit, loose regime & $+$ & $-2.25 \times 10^{-7}$ & -1.00 & 100\% & 0.009 & -- \\
\bottomrule
\end{tabular}

\end{table}

All three fail (Table~\ref{tab:mechanisms}), and none approaches the bar. The
largest $|t|$ among the primary tests is \mechMaxTCov. Each failure argues against a
further attempt. The H2 estimate is stable and negative on both discovery
(\mechHTwoBetaCov) and holdout (\mechHTwoHoldCov), which is the wrong sign, and it
never approaches significance. H3 is positive ($t=\mechHThreeTCov$) against a prior
of negative, so neither the magnitude nor the direction of the storage interaction
is present. H4's tight-regime effect is negative ($t=\mechHFourTightTCov$) and
smaller in magnitude than the loose-regime effect, which reverses both predictions.
On this evidence, cargo arrivals do not have a measurable effect in either storage
state.

The linear null of the previous subsection has thus been tested against the three
specific ways in which it could have been wrong, using pre-registered signs, a
corrected bar and an untouched holdout, and it stands.
